\documentclass[letterpaper,12pt]{article}
\usepackage[scale={0.8,0.9},centering,includeheadfoot,dvips, nohead]{geometry}
\usepackage[obeyspaces,spaces]{url}
\usepackage{amsmath}

\pagestyle{empty}






\newcommand{\hamcycle}{\textsc{HamCycle}\xspace}
\newcommand{\hampath}{\textsc{HamPath}\xspace}
\newcommand{\clique}{\textsc{Clique}\xspace}
\newcommand{\tsat}{\textsc{3SAT}\xspace}
\newcommand{\sss}{\textsc{Subset-Sum}\xspace}
\newcommand{\partition}{\textsc{Partition}\xspace}
\newcommand{\balance}{\textsc{Balance}\xspace}
\newcommand{\sws}{\textsc{SWS}\xspace}



\begin{document}

\noindent
 UCSC CSE201 Homework 2 \hfill Fall 2025 \\
Last update: 10/14/2025 \\
 \rule{\textwidth}{0.5pt}

%\input{../preamble}

\enlargethispage*{1.5cm}

%\noindent
% \rule{\textwidth}{0.5pt}




\begin{enumerate}


\item (10 pts) Assume that the hash table has a size $2^{10} = 1024$, and the universe $U := \{0, 1, 2, \ldots, 2^{1000}-1\}$. Consider a hash family $\mathcal{H}:= \{ h_b \; | \; b \in \{0, 1, 2, \ldots, 2^{10} -1\}\}$, where $h_b(x) = x + b \pmod{1024}$. Show that $\mathcal{H}$ is not a universal hash family. (Hint: You can do this by finding two elements/keys $u_1 \neq u_2$ in $U$ such that $\Pr_{h \sim \mathcal{H}} [ h(u_1) \neq h(u_2)] > 1/ 1024$.) 

\textbf{solution}

the problem lies with the fact that 1024 is not prime. 

consider \(u_1=0\) and \(u_2=1024\). when input into \(h_b\), they are guaranteed to result in the same number mod 1024 since \(u_1 \pmod {1024} = u_2 \pmod {1024} = 0\). 


\newpage
\item (10 pts) Assume that we have an input sequence of $n$ keys that should be mapped to a table of size $n$. Bob claims that his hash family is great by saying that for any table slot/bucket, the expected number of keys mapped to the slot is $O(1)$. Briefly explain why this is not a good way to measure the quality of hashing.


\textbf{solution}

the goal of a hash map should be to have the output space be smaller than the universe. furthermore, the process by which elements are mapped is not layed out so collisions of order \(O(n)\) are possible. 
\newpage


\item (20 pts) We are given an array of integers $A[1 \ldots n]$. We would like to know if there exists $i, j, k$ such that $A[i] = (A[j] + A[k])/2$. Solve this problem in $O(n^2)$ using hashing. 

nest two loops, sum over all \(i,j\) in n and divide by two. hash each value in the loops. this takes \(O(n^2)\) time. then loop over all n and hash each element. if there is a collision, assuming the hashing is collisionless, there exists \(i,j,k\) such that $A[i] = (A[j] + A[k])/2$

if we want to find \(jk\), we can store each hash value in a nxn matrix corresponding to the elements that lead to that value. checking the hashed value of \(A[i]\) against hte matrix for collisions should reveal the value of \(jk\) that leads to the equality. if there is a collision, assuming the hash is universal, there should be on the order of \(O(1)\) collisions so checking through each collision should take \(O(1)\) time.

\newpage



\item (30 pts) Consider the following variant of the Interval Selection Problem: We're given $n$ intervals, $I_i = (s_i, f_i)$, where $1 \leq i \leq n$. The length of interval $I_i$, denoted as $|I_i|$, is defined as $f_i - s_i$. Our goal is to choose a subset of mutually disjoint intervals with the maximum total length. That is, if $X$ is the set of indices of the chosen intervals, their total length will be $\sum_{i \in X} |I_i|$, and we want to maximize this value. We want to show that some natural greedy algorithms fail to find an optimum solution. Assume that all finish times and start times of the intervals are distinct. 

\begin{enumerate}
    \item Longest Interval First: In each iteration, choose a longest interval from the remaining intervals (and discard all intervals intersecting the chosen one). 
    
    \textbf{solution}
    
    consider n=3. the meetings span from \((0,4),(0,3),(3,6)\). the longest interval is the first one, but choosing the first means that no other meeting can be choosen. the optimal is choosing the second and third. the greedy results in \(t=4\), but the optimal is \(t=6\).

    
    \item Earliest Finishing Interval First: In each iteration, choose the interval ending the earliest from the remaining intervals.   
    
    \textbf{solution}

    consider n=2. meetings span from \((0,2),(0,3)\)

    greedy results in \(t=2\), picking only the first interval when the optimal is \(t=3\), picking only teh second interval.


    \item Earliest Starting Interval First: In each iteration, choose the interval starting  the earliest from the remaining intervals.        
    
    \textbf{solution}

    consider n=2. meetings span from \((0,2),(1,4)\)

    greedy results in \(t=2\), picking only the first interval when the optimal is \(t=3\), picking only the second interval.

\end{enumerate}

\newpage


\item (30 pts) There are two disjoint groups/sets of items, $A$ and $B$, i.e., $A \cap B = \emptyset$. Each item $i$ has a value $v_i$. We want to choose a subset of items with the maximum total value satisfying the following capacity constraints:
\begin{itemize}
    \item Choose at most $u_A$ items from set $A$.
    \item Choose at most $u_B$ items from set $B$.
    \item Choose at most $u_T$ items from the union of the two sets, $A \cup B$.
\end{itemize}

Consider the following greedy algorithm: Consider items in decreasing order of their value. In each iteration, Select the item if it doesn't violate any of the above capacity constraints. To show the optimality of this greedy algorithm, we will take two steps.

\begin{enumerate}
    \item Show that there exists an optimum solution that includes the first item our greedy algorithm selects. 
    
    \textbf{solution}

    % proceed by contradiction. let \(a\) be the highest value item in the current unselected \(u_a\) and let \(b\) be the highest value item in the current unselected \(u_b\), satisfying the first two requirements. suppose the item selected by our greedy algorithm was not in the optimum solution. at this step, item \(c\neq a,b\) was picked instead. swapping out item \(c\) for item \(o\in\{a,b\}\) satisfying \(\max(v(a),v(b))\), the item picked out by the algorithm, results in a higher value total, thus the item picked out by the algorithm must be part of the optimal solution. note that the swap is allowed because both \(a,b\) satisfy the first two requirements. 

    let \(g\) be the element selected by the greedy algorithm. if \(g\in OPT\), then we are done. suppose \(g\notin OPT\), if this is the case, then some other element \(c\) was choosen. but necessarily, \(v(c)\leq v(g)\). if they are equal, switching \(c\) with \(g\) results in the same value, or  \(v(c)< v(g)\) in which case swapping \(g\) with \(c\) results in a higher value, which is a contradiction. 
    % order the elements of \(A\cup B\) in descending value. let \(a\) be the first element that satisfies 



    \item Explain why this implies the optimality of the greedy algorithm. 
    
    if we run the algorithm \(\max(u_t,\lvert A\rvert + \lvert B\rvert)\) times, each time the algorithm runs the problem is converted into the subproblem of
    
    \begin{itemize}
    \item Choose at most $u_A-1$ items from set $A$.
    \item Choose at most $u_B$ items from set $B$.
    \item Choose at most $u_T-1$ items from the union of the two sets, $A \cup B$.
\end{itemize}

    or 
    \begin{itemize}
    \item Choose at most $u_A$ items from set $A$.
    \item Choose at most $u_B-1$ items from set $B$.
    \item Choose at most $u_T-1$ items from the union of the two sets, $A \cup B$.
\end{itemize}

    assuming we update the values of \(u_A,u_B,u_T\) at every step. 

    thus, we can define the original problem recursively until we hit the base case where either there are no more elements in \(A\) and \(B\), or \(u_T=0\). since at each step the item picked is optimal, the final solution is as well. 
    
    
\end{enumerate}






\end{enumerate}


\end{document}

